% Part: incompleteness
% Chapter: representability-in-q
% Section: beta-function

\documentclass[../../../include/open-logic-section]{subfiles}

\begin{document}

\olfileid[fr]{inc}{req}{bet}
\olsection{Le lemme de la fonction bêta}

Pour montrer que la récurrence primitive peut être réalisée à
l'aide de l'addition, de la multiplication et de~$\Char{=}$,
nous devons construire des fonctions qui manipulent des suites.
La tâche serait plus simple si nous disposions aussi de
l'exponentiation. Avec la récurrence primitive, nous pouvions
définir, par exemple, le « $n$-ième nombre premier » et adopter
un codage assez direct. Mais nous ne disposons pas ici de la
récurrence primitive : nous voulons précisément montrer qu'on
peut la réaliser par minimisation. Il faut donc trouver une
méthode plus ingénieuse.

\begin{lem}
\ollabel{lem:beta}
Il existe une fonction $\beta(d,i)$ telle que, pour toute suite
$a_0$, \dots,~$a_n$, il existe un nombre~$d$ vérifiant
$\beta(d,i)=a_i$ pour tout $i\le n$. De plus, $\beta$ peut être
définie à partir des fonctions de base par la seule composition
et la minimisation régulière.
\end{lem}

On peut voir $d$ comme un code de la suite
$\tuple{a_0,\dots,a_n}$ et $\beta(d,i)$ comme une fonction qui en
extrait le $i$-ième terme. Ce « codage » \emph{n'utilise pas} le
codage par puissances de nombres premiers déjà étudié. L'énoncé
du lemme est modeste : il ne dit ni que nous pouvons concaténer
des suites ou leur ajouter des termes, ni même que nous pouvons
\emph{calculer}~$d$ à partir de $a_0$, \dots,~$a_n$ avec des
fonctions définissables par composition et minimisation régulière.
Il affirme seulement l'existence d'une fonction de « décodage »
pour laquelle toute suite admet un code.

La notation $\beta$ vient de G\"odel. Rappelons que la difficulté
du lemme consiste à définir une fonction~$\beta$ adéquate avec
des moyens apparemment limités : la composition et la
minimisation, auxquelles s'ajoutent toutefois l'addition, la
multiplication et~$\Char{=}$. Il existe plusieurs démonstrations
du lemme, mais l'une des plus élégantes reste la méthode
originale de G\"odel. Elle repose sur un résultat d'arithmétique
appelé théorème de Sunzi, traditionnellement connu sous le nom
de « théorème des restes chinois ».

\begin{defn}
Deux entiers naturels $a$ et $b$ sont \emph{premiers entre eux} si
leur plus grand commun diviseur vaut~$1$; autrement dit, ils
n'ont aucun autre diviseur commun.
\end{defn}

\begin{defn}
Pour $c>0$, les entiers naturels $a$ et $b$ sont
\emph{congrus modulo~$c$},
ce que l'on note $a\equiv b\mod c$, si $c\mid(a-b)$, c'est-à-dire
s'ils ont le même reste dans la division par~$c$.
\end{defn}

Voici le théorème de Sunzi :
\begin{thm}
Supposons que $x_0$, \dots,~$x_n$ soient strictement positifs et
premiers entre eux deux
à deux. Soient $y_0$, \dots,~$y_n$ des entiers naturels quelconques.
Alors il existe un nombre~$z$ tel que
\begin{align*}
z & \equiv y_0 \mod x_0 \\
z & \equiv y_1 \mod x_1 \\
& \vdots  \\
z & \equiv y_n \mod x_n.
\end{align*}
\end{thm}

Voici comment nous utiliserons ce théorème. Si les nombres $x_0$,
\dots,~$x_n$ sont respectivement supérieurs à $y_0$,
\dots,~$y_n$, nous pouvons prendre $z$ comme code de la suite
$\tuple{y_0,\dots,y_n}$. Pour retrouver~$y_i$, il suffit de
diviser $z$ par~$x_i$ et de prendre le reste. Nous devons donc
choisir convenablement les nombres $x_0$, \dots,~$x_n$.

Deux observations nous y aideront. Étant donnés
$y_0$, \dots,~$y_n$, posons
\begin{align*}
j &= \max(n, y_0 + 1, \dots, y_n + 1), \\
m &= \lcm(1,\dots,j),
\end{align*}
où $m$ est le plus petit commun multiple de $1$, \dots,~$j$, et
posons
\begin{align*}
x_0 & = 1 + m \\
x_1 & = 1 + 2 \cdot m \\
x_2 & = 1 + 3 \cdot m \\
& \vdots  \\
x_n & = 1 + (n+1) \cdot m
\end{align*}
Deux propriétés en résultent :
\begin{enumerate}
\item\ollabel{rel-prime} $x_0,\dots,x_n$ sont premiers entre eux deux à deux.
\item\ollabel{less} Pour tout~$i$, on a $y_i<x_i$.
\end{enumerate}
Pour établir \olref{rel-prime}, supposons que $i\neq k$ et qu'un
nombre premier~$p$ divise $x_i$ et~$x_k$. Alors $p$ divise
$1+(i+1)m$ et $1+(k+1)m$, donc leur différence :
\[
(1 + (i+1)m) - (1+ (k+1)m) = (i-k) m.
\]
Comme $p$ divise $1+(i+1)m$, il ne peut aussi diviser~$m$ :
sinon, il diviserait~$1$. Puisqu'il divise $(i-k)m$ et qu'il
est premier, il divise donc $i-k$. Or $0<|i-k|\le n\le j$;
ainsi $p\le j$ et $p$ divise~$m$, le plus petit commun multiple
de $1$, \dots,~$j$. C'est une contradiction. Aucun nombre
premier ne divise donc deux des $x_i$. La propriété
\olref{less} est immédiate : $y_i<j\le m<x_i$.

Démontrons maintenant le lemme de la fonction~$\beta$. Rappelons
que nous pouvons employer la fonction nulle, le successeur,
l'addition, la multiplication, $\Char{=}$ et les projections,
ainsi que toute fonction obtenue à partir d'elles par composition
et minimisation appliquée à des fonctions régulières. Nous
pouvons également employer une relation si sa fonction
caractéristique se définit de cette manière. Comme précédemment,
on montre que ces relations sont closes par combinaisons
booléennes et par quantification bornée; par exemple :
\begin{align*}
\fn{not}(x) & \defis \Char{=}(x,0)\\
\bmin{x \leq z}{R(x,y)} & \defis \umin{x}{(R(x,y) \lor x = z)}\\
\bexists{x \leq z}{R(x,y)} & \defiff R(\bmin{x \leq z}{R(x,y)}, y)
\end{align*}
Nous pouvons alors montrer que les fonctions et relations
suivantes se définissent elles aussi sans récurrence primitive :
\begin{enumerate}
\item la fonction de couplage
  $J(x,y)=\frac{1}{2}[(x+y)(x+y+1)]+x$;
\item les fonctions de projection
\begin{align*}
K(z) & = \bmin{x \leq z}{\bexists{y \leq z}{z = J(x,y)}},\\
L(z) & = \bmin{y \leq z}{\bexists{x \leq z}{z = J(x,y)}};
\end{align*}
\item la relation d'ordre strict $x<y$;
\item la relation de divisibilité $x\mid y$;
\item la fonction $\fn{rem}(x,y)$, qui donne le reste de la division
  de~$y$ par~$x$.
\end{enumerate}
Nous posons $\fn{rem}(0,y)=0$ pour rendre cette fonction totale;
les diviseurs utilisés ci-dessous sont tous positifs.
Définissons
\begin{align*}
\beta^*(d_0,d_1,i) & = \fn{rem}(1+(i+1) d_1,d_0) \text{ et}\\
\beta(d,i) & = \beta^*(K(d),L(d),i).
\end{align*}
C'est la fonction cherchée. Étant donnés
$a_0,\dots,a_n$ comme plus haut, posons
\[
j = \max(n,a_0+1,\dots,a_n+1),
\]
puis $d_1=\lcm(1,\dots,j)$. D'après \olref{rel-prime}, les
nombres $1+d_1$, $1+2d_1$, \dots, $1+(n+1)d_1$ sont premiers
entre eux deux à deux et, d'après \olref{less}, ils sont
respectivement supérieurs à $a_0$, \dots, $a_n$. Le théorème de
Sunzi fournit donc un nombre~$d_0$ tel que, pour chaque~$i$,
\[
d_0 \equiv a_i \mod (1+(i+1)d_1)
\]
et, puisque $d_1>a_i$,
\[
a_i = \fn{rem}(1+(i+1)d_1,d_0).
\]
Soit $d=J(d_0,d_1)$. Pour tout $i\le n$, nous avons
\begin{align*}
\beta(d,i) & =  \beta^*(d_0,d_1,i) \\
& =  \fn{rem}(1+(i+1) d_1,d_0) \\
& =  a_i
\end{align*}
comme il fallait le montrer. La preuve du lemme de la fonction
$\beta$ est achevée.

\begin{prob}
  Montrez que les relations $x<y$ et $x\mid y$, ainsi que la
  fonction~$\fn{rem}(x,y)$, se définissent sans récurrence
  primitive. Vous pouvez employer la fonction nulle, le
  successeur, l'addition, la multiplication, $\Char{=}$, les
  projections, ainsi que la minimisation et la quantification
  bornées.
\end{prob}

\begin{explain}
Dans la définition de~$\beta$, on calcule le module à partir
de $L(d)$ et de la position~$i$, puis on lit le reste de
$K(d)$ modulo ce nombre. Le choix d'un code pour une suite
donnée intervient seulement dans la preuve d'existence. Pour
l'application suivante, il suffira que chaque suite finie de
valeurs admette au moins un code.
\end{explain}

\end{document}
